% 第4章 実験。oshifit_spec.md 第4〜6節にもとづく下書き。
\chapter{実験}
\section{目的}
4週間の使用を通じて、提案の精度と、幸福度・ストレスの変化を調べる。

\section{参加者}
\todo{人数、募集方法、選択・除外の基準、謝礼}

\section{手続き}
参加者は、説明を読んで同意したのち、最初の調査に答え、アバターを作る。
分析の結果にもとづいて群A・B・Cのいずれかになり、群B・Cの参加者は約4週間、4〜5日ごとに6回の振り返りを行う。
28日目以降に最後の調査に答える。
実験は \todo{倫理審査の承認日} から2027年3月31日までの研究期間の中で行う。

\section{測定}
幸福度は WHO-5（0〜100に換算）、ストレスは PSS-J（0〜56）で測る。
副次的に、主観的幸福感（SHS）、推し疲れ、5つの要素を用いる。

\section{分析}
\subsection{提案の精度}
示した見込みと28日目の実測を比べ、平均絶対誤差と、見込みの幅に実測が入った割合を求める。
\subsection{心理的な効果}
群ごとに開始時から28日目への変化を対応のある検定で調べ、効果量を求める。
群Bと群Cの比較には、28日目の値を目的変数、開始時の値と性格を共変量とする共分散分析を用いる。
群は本人の選択によるため、傾向スコアによる重みづけを感度分析として行う。
\todo{有意水準、多重比較、欠測の扱い}

\section{倫理的配慮}
本研究は \todo{倫理審査委員会の名称と承認番号} の承認を得て行った。
